\section{Conclusion}
\label{sec:conclusion}

We began with a widely held assumption: contrastive SSL amplifies spurious correlations due to augmentation invariance. Our experiments demonstrate the opposite.

Supervised ERM encodes spurious background features significantly more strongly than contrastive MoCo-v3 on Waterbirds (ratio diff = 0.025, Cohen's $d = 5.68$, $p < 0.0001$), while ERM $\approx$ DINO ($p = 1.0$). The pretraining paradigm effect is highly significant across both datasets (Waterbirds ANOVA $F = 35.99$, $p = 2.42 \times 10^{-7}$; CelebA $F = 5.51$, $p = 0.009$). The paradigm ranking reverses across datasets, establishing that spurious attribute type interacts with pretraining objective in ways not captured by existing theory.

\textbf{Contributions summarized:} (1) First controlled 4-paradigm comparison of spurious feature encoding in frozen ResNet-50 representations. (2) Supervised label correlation identified as the dominant driver --- label-agnostic pretraining suppresses spurious encoding. (3) Spurious-attribute-type $\times$ pretraining-objective interaction established via cross-dataset ranking reversal. (4) Background-replacement augmentation verified as a functional causal lever.

\textbf{Future directions} grounded in our results: (i) MAE comparison to distinguish label-agnostic training from contrastive augmentation effects; (ii) downstream WGA evaluation via DFR on MoCo-v3 vs ERM features; (iii) ViT backbone and additional spurious attribute type comparison to map the full interaction matrix.

The pretraining objective determines how much spurious content enters the representation --- and label-agnostic pretraining provides partial spurious feature suppression as a free byproduct of not seeing class labels. Practitioners should factor this into backbone selection when spurious correlation robustness matters.
